\uuid{RebY}
\titre{Application linéaire}
\auteur{hanani ; mzari}
\organisation{exo7}
\licence{CC BY-NC-SA 4.0}
\datecreate{2020-01-20}
\niveau{L1}
\module{Algèbre}
\chapitre{Application linéaire}
\sousChapitre{Définition}
\difficulte{1}
\theme{Applications linéaires, Application linéaire}
\exo7qcmid{lille-2.applications-lineaires.001}

\contenu{
\begin{question}
\qclassification{107.01}
\qtags{motcle=Application linéaire}
On considère les deux applications suivantes :
$$\begin{array}{rccc}f:&\Rr&\to&\Rr\\
& x&\to & \sin x \end{array} \quad \mbox{et} \quad \begin{array}{rccc}g:&\Rr^2&\to&\Rr^2\\
& (x,y)&\to &(y,x). \end{array}$$
Quelles sont les assertions vraies ?
\begin{answers}
\good{$f(0)=0$}.
\bad{$f$ est une application linéaire}.
\bad{$g(x,y) =  g(y,x)$, pour tout $(x,y) \in \Rr^2$}.
\good{$g$ est une application linéaire}.
\end{answers}
\begin{explanations} $f$ n'est pas linéaire car $f(\pi)=f(\frac{\pi}{2}+ \frac{\pi}{2})=0$
et $f(\frac{\pi}{2})+ f(\frac{\pi}{2}) =2$.
On vérifie que $g$ est linéaire.
\end{explanations}
\end{question}
}
